\section*{Capítulo 1: Problema 1}

\textbf{Ejercicios:}

En cada caso, determine cuáles vectores localizados $\overrightarrow{PQ}$ y $\overrightarrow{AB}$ son equivalentes.

\begin{enumerate}
    \item $P = (1, -1)$, $Q = (4, 3)$, $A = (-1, 5)$, $B = (5, 2)$.
    
    \textbf{Solución:}
    \begin{align*}
        \overrightarrow{PQ} &= Q - P = (4 - 1, 3 - (-1)) = (3, 4), \\
        \overrightarrow{AB} &= B - A = (5 - (-1), 2 - 5) = (6, -3).
    \end{align*}
    Los vectores $\overrightarrow{PQ}$ y $\overrightarrow{AB}$ no son equivalentes.

    \item $P = (1, 4)$, $Q = (-3, 5)$, $A = (5, 7)$, $B = (1, 8)$.
    
    \textbf{Solución:}
    \begin{align*}
        \overrightarrow{PQ} &= Q - P = (-3 - 1, 5 - 4) = (-4, 1), \\
        \overrightarrow{AB} &= B - A = (1 - 5, 8 - 7) = (-4, 1).
    \end{align*}
    Los vectores $\overrightarrow{PQ}$ y $\overrightarrow{AB}$ son equivalentes.

    \item $P = (1, -1, 5)$, $Q = (-2, 3, -4)$, $A = (3, 1, 1)$, $B = (0, 5, 10)$.
    
    \textbf{Solución:}
    \begin{align*}
        \overrightarrow{PQ} &= Q - P = (-2 - 1, 3 - (-1), -4 - 5) = (-3, 4, -9), \\
        \overrightarrow{AB} &= B - A = (0 - 3, 5 - 1, 10 - 1) = (-3, 4, 9).
    \end{align*}
    Los vectores $\overrightarrow{PQ}$ y $\overrightarrow{AB}$ no son equivalentes.

    \item $P = (2, 3, -4)$, $Q = (-1, 3, 5)$, $A = (-2, 3, -1)$, $B = (-5, 3, 8)$.
    
    \textbf{Solución:}
    \begin{align*}
        \overrightarrow{PQ} &= Q - P = (-1 - 2, 3 - 3, 5 - (-4)) = (-3, 0, 9), \\
        \overrightarrow{AB} &= B - A = (-5 - (-2), 3 - 3, 8 - (-1)) = (-3, 0, 9).
    \end{align*}
    Los vectores $\overrightarrow{PQ}$ y $\overrightarrow{AB}$ son equivalentes.
\end{enumerate}

En cada caso, determine cuáles vectores localizados $\overrightarrow{PQ}$ y $\overrightarrow{AB}$ son paralelos.

\begin{enumerate}
    \setcounter{enumi}{4}
    \item $P = (1, -1)$, $Q = (4, 3)$, $A = (-1, 5)$, $B = (7, 1)$.
    
    \textbf{Solución:}
    \begin{align*}
        \overrightarrow{PQ} &= Q - P = (4 - 1, 3 - (-1)) = (3, 4), \\
        \overrightarrow{AB} &= B - A = (7 - (-1), 1 - 5) = (8, -4).
    \end{align*}
    Los vectores $\overrightarrow{PQ}$ y $\overrightarrow{AB}$ no son paralelos.

    \item $P = (1, 4)$, $Q = (-3, 5)$, $A = (5, 7)$, $B = (9, 6)$.
    
    \textbf{Solución:}
    \begin{align*}
        \overrightarrow{PQ} &= Q - P = (-3 - 1, 5 - 4) = (-4, 1), \\
        \overrightarrow{AB} &= B - A = (9 - 5, 6 - 7) = (4, -1).
    \end{align*}
    Los vectores $\overrightarrow{PQ}$ y $\overrightarrow{AB}$ son paralelos.

    \item $P = (1, -1, 5)$, $Q = (-2, 3, -4)$, $A = (3, 1, 1)$, $B = (-3, 9, -17)$.
    
    \textbf{Solución:}
    \begin{align*}
        \overrightarrow{PQ} &= Q - P = (-2 - 1, 3 - (-1), -4 - 5) = (-3, 4, -9), \\
        \overrightarrow{AB} &= B - A = (-3 - 3, 9 - 1, -17 - 1) = (-6, 8, -18).
    \end{align*}
    Como $\overrightarrow{AB} = 2 \cdot \overrightarrow{PQ}$, los vectores $\overrightarrow{PQ}$ y $\overrightarrow{AB}$ son paralelos.

    \item $P = (2, 3, -4)$, $Q = (-1, 3, 5)$, $A = (-2, 3, -1)$, $B = (-11, 3, -28)$.
    
    \textbf{Solución:}
    \begin{align*}
        \overrightarrow{PQ} &= Q - P = (-1 - 2, 3 - 3, 5 - (-4)) = (-3, 0, 9), \\
        \overrightarrow{AB} &= B - A = (-11 - (-2), 3 - 3, -28 - (-1)) = (-9, 0, -27).
    \end{align*}
    Como $\overrightarrow{AB} = 3 \cdot \overrightarrow{PQ}$, los vectores $\overrightarrow{PQ}$ y $\overrightarrow{AB}$ son paralelos.
\end{enumerate}

\textbf{Ejercicio 9:}

Dibuje los vectores localizados de los ejercicios 1, 2, 5 y 6 en una hoja de papel para ilustrar estos ejercicios. También dibuje los vectores localizados $\overrightarrow{QP}$ y $\overrightarrow{BA}$. Dibuje los puntos $Q - P$, $B - A$, $P - Q$ y $A - B$.